\section{Introduction}

\subsection{Context}

Modern organisations are drowning in data---structured, unstructured, and
multimodal. The competitive challenge, however, is no longer storing data but
extracting the precise insight needed at the right moment. In sectors such as
legal discovery, fraud detection, intelligence analysis, healthcare diagnostics,
and regulatory compliance, the ability to locate a single rare but critical fact
within terabytes of irrelevant or misleading material---the proverbial
\emph{needle in a haystack}---has become an operational imperative. Foundation
models such as IBM Granite, designed for enterprise-grade performance, offer a
promising route to retrieval systems that understand context, intent, and
domain-specific nuance rather than merely matching keywords.

\subsection{The Problem}

Traditional keyword search and basic retrieval tools often fail to surface rare,
relevant information, because relevance is semantic rather than lexical: the
right passage may share few words with the query. Dense (semantic) retrieval and
the retrieval-augmented generation (RAG) introduced by \textcite{lewis2020rag}
promise to close this gap, but a practitioner choosing a system needs evidence: does a
Granite-based retrieval system actually find relevant information more accurately
than the established alternatives, and can its answers be trusted?

\subsection{Gap}

Two gaps motivate this project. First, while RAG and dense retrieval are widely
discussed, there is little rigorous, like-for-like evidence comparing an
\emph{enterprise-grade Granite} retrieval system against standard baselines on
the established information-retrieval benchmarks that come with human relevance
judgments. Second, the brief explicitly calls for \emph{explainability and trust
in retrieved outputs}---yet most retrieval demonstrations return answers without
attributing them to their sources, leaving users unable to verify a result.

\subsection{Proposed Solution}

This project designs and evaluates a Granite-powered retrieval system for precise
information discovery. The system combines a dense semantic retriever with a RAG
generation layer and an explainability component that attributes each answer to
the passages supporting it. Its retrieval quality is measured on standard IR
benchmarks (BEIR \parencite{thakur2021beir}, MS~MARCO \parencite{nguyen2016msmarco},
Natural Questions \parencite{kwiatkowski2019nq}) using precision@k, recall@k,
nDCG, and MRR, and compared head-to-head against a classical BM25 baseline
\parencite{robertson2009bm25} and an open-source dense retriever
\parencite{karpukhin2020dpr}. As a secondary diagnostic, a long-context
``needle'' stress test examines how retrieval degrades when the relevant passage
is buried deep within a long document---the \emph{Lost-in-the-Middle} effect
\parencite{liu2024lost}. The result is concrete, decision-relevant evidence about
when and where a Granite-based system improves enterprise information discovery.
